\documentclass[twocolumn]{aastex631}
\begin{document}
\title{A residual reionization ionizing-photon-budget shortfall for star-forming galaxies at z\~{}8 (beta systematic corner)}
\author{NebulaMind Lab (autonomous pipeline)}
\affiliation{NebulaMind Lab, \url{https://lab.nebulamind.net}}
\begin{abstract}
Reionization ionizing-photon-budget reconciliation at z\~{}8: star-forming galaxies require f\_esc=0.178 (+0.179/-0.091) to close the budget (Madau-Dickinson SFRD, log xi\_ion=25.5+/-0.15, clumping C=2-5, JWST-SFRD tail), versus indirect-proxy-inferred f\_esc=0.050 (+0.076/-0.030) from LzLCS O32/beta calibrations. Median delta(required-inferred)=+0.115 dex-frac (16-84\%: +0.009 to +0.294); 86\% of the systematic MC shows a shortfall. Conclusion: a genuine SHORTFALL remains, and the sign holds under both O32 and beta calibrations. The result is bounded by the xi\_ion x clumping x proxy-calibration systematic, not by statistics. Generated autonomously from public data (jwst) via the NebulaMind Lab runner: a bounded, reproducible, descriptive study.
\end{abstract}
\section{Introduction}
Reconciling the reionization ionizing-photon-budget has been a longstanding challenge in understanding cosmic history. Recent studies have highlighted potential shortfalls in the photon budget at high redshifts [Muñoz2024], suggesting that current models may not fully account for the required ionizing photons to drive reionization. This issue is further complicated by uncertainties in key parameters such as escape fraction (f\_esc) and ionizing efficiency (xi\_ion). Previous works have emphasized the importance of accurately calibrating these values to avoid underestimating the photon budget [Davies2021, Park2022]. Building on this foundation, our study aims to reassess the reionization photon budget using a literature-anchored approach.
\section{Data and method}
To address this challenge, we adopt the Madau \& Dickinson (2014) cosmic star formation rate density (SFRD) as our starting point. We then incorporate published values for xi\_ion and f\_esc proxy calibrations from LzLCS [Chisholm+22, Flury+22] and Simmonds+24. Notably, we do not rely on any new observational data or survey catalogs; instead, we focus on reconciling existing literature values to identify potential discrepancies in the photon budget.
\section{Result}
Our analysis reveals a significant shortfall in the reionization ionizing-photon-budget at z\~{}8. Specifically, star-forming galaxies require an escape fraction of f\_esc=0.178 (+0.179/-0.091) to close the budget, assuming the Madau-Dickinson SFRD, log xi\_ion=25.5+/-0.15, and clumping factor C=2-5. However, indirect-proxy-inferred values from LzLCS O32/beta calibrations yield a lower f\_esc=0.050 (+0.076/-0.030). This discrepancy results in a median delta of +0.115 dex-frac (16-84\%: +0.009 to +0.294), with 86\% of systematic Monte Carlo simulations showing a shortfall.
\begin{figure}\centering\includegraphics[width=\columnwidth]{result.png}\caption{A residual reionization ionizing-photon-budget shortfall for star-forming galaxies at z\~{}8 (beta systematic corner)}\end{figure}
\section{Caveats}
It is essential to acknowledge the limitations of our approach. Our analysis relies on automated, single-selection, and uncalibrated measurements, which may introduce biases or uncertainties not fully accounted for in our study. Additionally, the use of published proxy calibrations can be sensitive to variations in underlying assumptions and model parameters. Furthermore, our reliance on literature values means that we inherit any systematic errors present in those studies. Therefore, while our results highlight a potential shortfall in the reionization photon budget, further investigation using more comprehensive datasets and refined calibration methods is necessary to confirm these findings.
\section*{References}
\footnotesize
[Muoz2024] Muñoz et al. (2024). Reionization after JWST: a photon budget crisis?. bibcode 2024MNRAS.535L..37M \par [Davies2021] Davies et al. (2021). The Predicament of Absorption-dominated Reionization: Increased Demands on Ionizing Sources. bibcode 2021ApJ...918L..35D \par [Park2022] Park et al. (2022). Calibrating excursion set reionization models to approximately conserve ionizing photons. bibcode 2022MNRAS.517..192P \par [Duncan2015] Duncan et al. (2015). Powering reionization: assessing the galaxy ionizing photon budget at z < 10. bibcode 2015MNRAS.451.2030D \par [Madau2017] Madau (2017). Cosmic Reionization after Planck and before JWST: An Analytic Approach. bibcode 2017ApJ...851...50M

\end{document}
